% MyGrowth - project portfolio
% Compile with: pdflatex portfolio.tex  (or upload to Overleaf)
% Every screenshot is a placeholder. Replace \shot{...} with
% \includegraphics[width=2.4cm]{screens/<file>.png} once you have captures.
% To add another project, copy the \project{...} block and its sections.

\documentclass[11pt,a4paper]{article}

\usepackage[margin=2cm]{geometry}
\usepackage[T1]{fontenc}
\usepackage{lmodern}
\usepackage{microtype}
\usepackage{graphicx}
\usepackage[dvipsnames]{xcolor}
\usepackage{tikz}
\usepackage{enumitem}
\usepackage{titlesec}
\usepackage{fancyhdr}
\usepackage[hidelinks]{hyperref}

% ---------- Palette ----------
\definecolor{teal}{HTML}{247F7B}
\definecolor{mind}{HTML}{5B5FC7}
\definecolor{fuel}{HTML}{3A8A4F}
\definecolor{move}{HTML}{C15F2C}
\definecolor{night}{HTML}{4B4FA0}
\definecolor{trail}{HTML}{DD8C43}
\definecolor{ink}{HTML}{25393B}
\definecolor{muted}{HTML}{6B7F80}
\definecolor{soft}{HTML}{EEF6F4}
% Mind branch colours (from the app)
\definecolor{selfacc}{HTML}{5E8C61}
\definecolor{joy}{HTML}{B7791F}
\definecolor{purpose}{HTML}{3D7F8F}
\definecolor{livelife}{HTML}{2B8273}
\definecolor{focus}{HTML}{3F6FC4}
\definecolor{standup}{HTML}{C2603A}
\definecolor{aware}{HTML}{7E5DB5}
\definecolor{control}{HTML}{B5527A}
\definecolor{social}{HTML}{C9802A}

\color{ink}
\setlength{\parindent}{0pt}
\setlength{\parskip}{5pt}
\setlist{itemsep=1pt, topsep=3pt, leftmargin=1.1em}

% ---------- Headings ----------
\titleformat{\section}{\large\bfseries\color{teal}}{}{0pt}{}[\vspace{-5pt}{\color{teal!30}\rule{\linewidth}{0.6pt}}]
\titlespacing*{\section}{0pt}{12pt}{6pt}

\pagestyle{fancy}
\fancyhf{}
\renewcommand{\headrulewidth}{0pt}
\fancyfoot[L]{\small\color{muted}[Your Name] \textperiodcentered{} Portfolio}
\fancyfoot[R]{\small\color{muted}\thepage}

% ---------- Helpers ----------
% Small phone-shaped screenshot placeholder: \shot{Caption}
\newcommand{\shot}[1]{%
  \begin{tikzpicture}
    \draw[muted!60, line width=1pt, rounded corners=7pt, fill=soft] (0,0) rectangle (2.4,4.6);
    \draw[muted!50, line width=0.8pt] (0.95,4.43) -- (1.45,4.43);
    \draw[muted!50, dashed, rounded corners=3pt] (0.17,0.22) rectangle (2.23,4.27);
    \node[align=center, text width=1.9cm, font=\tiny\color{muted}] at (1.2,2.25)
      {\textbf{SCREENSHOT}\\[2pt]#1};
  \end{tikzpicture}%
}

% Real screenshot with a rounded frame and caption: \screen{file}{Caption}
\newcommand{\screen}[2]{%
  \begin{minipage}[t]{2.4cm}\centering
    \begin{tikzpicture}
      \node[inner sep=0pt](img){\includegraphics[width=2.4cm,height=4.6cm,keepaspectratio]{#1}};
      \draw[muted!50, line width=0.8pt, rounded corners=4pt]
        (img.south west) rectangle (img.north east);
    \end{tikzpicture}\par\vspace{2pt}
    {\tiny\color{muted}#2}
  \end{minipage}%
}

% Coloured tag: \pill{colour}{Text}
\newcommand{\pill}[2]{%
  \tikz[baseline=(t.base)]\node[fill=#1!15, text=#1!70!black, rounded corners=5pt,
    inner xsep=5pt, inner ysep=1.5pt, font=\scriptsize\bfseries](t){#2};%
}

% Mind branch row: \branch{colour}{Time / theme}{Branch}{What you do there}
\newcommand{\branch}[4]{%
  \pill{#1}{#2}\,{\bfseries\color{#1!80!black}#3}\par
  {\small #4}%
}

% Project banner: \project{colour}{Title}{One-line summary}
\newcommand{\project}[3]{%
  \par\noindent
  \tikz\node[fill=#1, rounded corners=10pt, inner sep=14pt,
             text width=\dimexpr\linewidth-28pt\relax, align=left, text=white]{%
    {\scriptsize\bfseries PROJECT}\par\vspace{2pt}
    {\LARGE\bfseries #2}\par\vspace{3pt}
    {\small #3}};\par\vspace{6pt}
}

\begin{document}

% ======================================================================
% PROJECT: MyGrowth
% ======================================================================
\project{teal}{MyGrowth}{A mobile wellbeing companion for students and young adults
  \textperiodcentered{} Lead Developer \textperiodcentered{} 2026\\
  \href{https://github.com/pskpsk009/BetterDays}{\texttt{github.com/pskpsk009/BetterDays}}}

MyGrowth turns small wellbeing habits into short, hands-on activities. Users
tap the head, body, or legs of an interactive figure on the home screen to
open \textbf{Mind}, \textbf{Body}, or the \textbf{Movement Trail}. Every
activity is tap-, drag-, or sensor-driven rather than journal-style, and
nothing is scored.

\pill{teal}{Expo SDK 57} \pill{teal}{React Native} \pill{teal}{TypeScript}
\pill{teal}{Expo Router} \pill{mind}{Reanimated} \pill{mind}{react-native-svg}
\pill{trail}{expo-location} \pill{trail}{expo-sensors}
\pill{trail}{react-native-maps} \pill{fuel}{Supabase}

% ----------------------------------------------------------------------
\section{Home \& Navigation}

\begin{center}
\screen{MentalHome.jpg}{Home: interactive body figure}\hfill
\screen{Mind.jpg}{Mind tab: branch list}\hfill
\screen{SelfAccept.jpg}{Inside a branch: Self-Acceptance}\hfill
\screen{LetterToPast.jpg}{An activity: Letter to My Past Self}
\end{center}

\begin{minipage}[t]{0.48\linewidth}
\pill{teal}{Home}
\begin{itemize}
  \item An \textbf{SVG body figure} with three tap zones:
        \textbf{Head} opens Mind, \textbf{Body} opens Body Functions, and
        \textbf{Legs} open the Movement Trail.
  \item Dashed callout labels point to each zone. Tapping one highlights it
        in its colour, dims the others, and shows a ``Selected area'' toast
        for a moment before opening that section.
  \item A softly \textbf{breathing halo} pulses behind the figure. It stays
        still when the phone's reduce-motion setting is on.
\end{itemize}
\end{minipage}\hfill
\begin{minipage}[t]{0.48\linewidth}
\pill{teal}{Navigation}
\begin{itemize}
  \item \textbf{Expo Router} with the phone's own \textbf{native bottom tabs}:
        Home, Mind, Body, Trail, and Progress, with platform icons on iOS
        and Android.
  \item Mind and Body share one \textbf{branch explorer}: section overview
        $\rightarrow$ branch $\rightarrow$ activity, with a back button at each
        level. Mind adds a Mental/Emotional switch at the top.
  \item Every branch has its own colour and icon, which carry through to
        its activities so users always know where they are.
\end{itemize}
\end{minipage}

\newpage
% ----------------------------------------------------------------------
\section{Mind: Mental \& Emotional Wellness}

\begin{center}
\screen{MentalHealth.jpg}{Mental Wellness overview}\hfill
\screen{MentalBranch.jpg}{A Mental Wellness branch}\hfill
\screen{EmotionalHealth.jpg}{Emotional Wellness overview}\hfill
\screen{EmotinalBranch.jpg}{An Emotional Wellness branch}
\end{center}

\begin{minipage}[t]{0.48\linewidth}
\pill{mind}{Mental Wellness}\par\vspace{2pt}
{\small\color{muted}Grow from the past, enjoy today, look ahead with focus.}
\begin{itemize}[leftmargin=0pt, label={}]
  \item \branch{selfacc}{Past}{Self-Acceptance}{Make peace with the past:
        lessons not labels, and a garden of your strengths.}
  \item \branch{livelife}{Present}{Live the Life}{Notice the moment with the
        Breathing Bubble, 5-Senses Grounding, and Watch Your Thoughts.}
  \item \branch{joy}{Today}{Discover Your Joy}{Explore hobbies and pick one
        small, new thing to try today.}
  \item \branch{purpose}{Future}{Purpose in Life}{Choose your values and one
        manageable step toward a goal.}
  \item \branch{focus}{Every day}{Find Your Focus}{Train attention with Focus
        Sprint and the Target Tap game.}
\end{itemize}
\end{minipage}\hfill
\begin{minipage}[t]{0.48\linewidth}
\pill{night}{Emotional Wellness}\par\vspace{2pt}
{\small\color{muted}Understand your feelings, bounce back, connect with others.}
\begin{itemize}[leftmargin=0pt, label={}]
  \item \branch{standup}{Bounce back}{Stand Up Again}{Turn a setback into one
        next step, and find a Plan B when plans fall through.}
  \item \branch{aware}{Understand}{Self-Awareness}{Name a feeling with the
        Emotion Explorer and sort what you can and can't control.}
  \item \branch{control}{Steady}{Self-Control}{Stop, notice, breathe, then
        choose how to respond.}
  \item \branch{social}{Connect}{Social Creatures}{Practise conversation
        openers, other points of view, and kind boundaries.}
\end{itemize}
\end{minipage}

% ----------------------------------------------------------------------
\section{Body Functions}

\begin{center}
\screen{Hydration.jpg}{Hydration Check-in}\hfill
\screen{BalancedDiet.jpg}{Build a Balanced Plate}\hfill
\screen{DailyStretch.jpg}{Daily Stretch}\hfill
\screen{SleepSchedule.jpg}{Sleep Schedule Guidance}
\end{center}

\begin{itemize}
  \item \pill{fuel}{Fuel} \textbf{Hydration Check-in}: tap a glass per drink,
        with an adjustable glass size (150--500\,ml) and a daily target you
        set yourself (1--4\,L).
  \item \pill{fuel}{Fuel} \textbf{Build a Balanced Plate}: long-press and drag
        food groups onto a five-slice SVG plate. The slice under your finger
        is found from its angle to the centre, and the feedback describes
        rather than judges.
  \item \pill{move}{Movement} \textbf{Daily Stretch} and \textbf{Desk Break}:
        guided move cards with countdown timers (3, 5, or 10-minute stretch
        routines) and a reminder to stop if anything hurts.
  \item \pill{night}{Sleep} \textbf{Sleep Schedule Guidance}: pick current
        and preferred bedtimes on a 24-hour picker and get a gradual
        night-by-night plan in 15 or 30-minute steps.
\end{itemize}

\newpage
% ----------------------------------------------------------------------
\section{Movement Trail}

\begin{center}
\screen{Outdoor.jpg}{Outdoor GPS trail}\hspace{1cm}
\screen{Indoor.jpg}{Indoor 2D drawing}\hspace{1cm}
\screen{ViewTrail.jpg}{View Trails gallery}
\end{center}

\begin{itemize}
  \item \pill{trail}{Outdoor} Live GPS from \texttt{expo-location} drawn on
        \texttt{react-native-maps}, with haversine distance. Pause and resume
        start a new segment, so ground covered while paused isn't counted.
  \item \pill{trail}{Indoor} For places without good GPS: accelerometer peaks
        detect steps and the magnetometer gives each step a direction. The
        path is drawn on an SVG canvas that scales to fit and smooths the line.
  \item \pill{trail}{View Trails} Stopping offers \emph{Save image} or
        \emph{Discard}. Saved trails are uploaded as SVGs to Supabase Storage
        and appear in a shared gallery, newest first.
\end{itemize}

% ----------------------------------------------------------------------
\section{What I learned \& next steps}

\begin{itemize}
  \item \textbf{Phone sensors are noisy.} Raw accelerometer data counts too
        many steps, so a step needs the signal to settle and a short gap first.
        Compass drift indoors is why the indoor trail is a drawing, not a map.
  \item \textbf{Layout happens after first render.} SVG backgrounds sized in
        percentages didn't grow when text wrapped; measuring the real layout
        fixed overflowing headers.
  \item \textbf{Keep logic out of screens.} Timing, game rules, plate
        feedback, and sleep scheduling live in plain TypeScript modules, which
        made them easier to change.
  \item \textbf{Next:} background trail tracking, real saved progress (with
        consent, export, and deletion), automated tests for the logic
        modules, and moderation for the shared trails gallery.
\end{itemize}

{\footnotesize\color{muted} MyGrowth is a prototype for everyday habit
support. It does not diagnose, treat, or replace professional care.}

% ======================================================================
% NEXT PROJECT: copy the \project{...} line and sections above, e.g.
% \newpage
% \project{move}{Project Two}{Short summary \textperiodcentered{} [Role] \textperiodcentered{} [Year]}
% ======================================================================

\end{document}
